%% ============================================================
%% NEW SECTION: The Octonionic Extension
%% Insert AFTER Section 7.4 "The Gauge-Gravity Separation" (~line 1081)
%% and BEFORE Section 8 "Discussion and Conclusions" (~line 1084)
%% ============================================================

\section{The Octonionic Extension: Three Generations from Geometry}
\label{sec:4-octonionic}

The Chamseddine--Connes--Marcolli (CCM) construction of Sections~\ref{sec:4-brane-triple}--\ref{sec:4-gauge-gravity} takes the number of fermion generations $N_g$ as a free parameter: the finite algebra $\mathcal{A}_F = \mathbb{C} \oplus \mathbb{H} \oplus M_3(\mathbb{C})$ encodes the gauge group $G_{\mathrm{SM}} = \mathrm{SU}(3)_c \times \mathrm{SU}(2)_L \times \mathrm{U}(1)_Y$ but says nothing about $N_g$. One simply declares $\mathcal{H}_F = \mathbb{C}^{32 N_g}$ and sets $N_g = 3$ by hand. In this section we show that extending the finite algebra from the direct sum $\mathbb{C} \oplus \mathbb{H} \oplus M_3(\mathbb{C})$ to the complexified \emph{Dixon algebra}
\begin{equation}
  \mathcal{T}_{\mathbb{C}} = \mathbb{C}_{\mathbb{C}} \otimes_{\mathbb{C}} \mathbb{H}_{\mathbb{C}} \otimes_{\mathbb{C}} \mathbb{O}_{\mathbb{C}}
  \label{eq:4-74}
\end{equation}
derives $N_g = 3$ as an algebraic rigidity result, while preserving the gauge group, the spectral action coefficients $(C^2, E_4, R^2) = (-18, +11, 0)$, the conformal coupling $\xi = 1/6$, and all bulk physics.

The construction rests on four independent algebraic theorems (Hurwitz, Albert, triality uniqueness, and the Fano plane structure), each forcing $N_g \leq 3$ through a different mechanism. The resulting spectral triple $(\mathcal{T}_{\mathbb{C}}, \mathcal{H}_{\mathrm{oct}}, D_{\mathrm{oct}}, J_{\mathrm{oct}}, \gamma_{\mathrm{oct}})$ satisfies all seven NCG axioms---with the first-order condition holding in the Boyle--Farnsworth modified sense for alternative algebras~\cite{ch4:BoyleFarnsworth2020}. We verify orientability by two independent proofs (envelope surjectivity and explicit Hochschild cycle) and establish Poincar\'{e} duality through the associative envelope $A_{\mathrm{env}}(\mathbb{O}) = M_8(\mathbb{R})$.

The division of labour between the algebraic (octonionic) and geometric (5D warp) sectors is clean: the octonions determine the \emph{structure}---$N_g = 3$, the gauge group, the $S_3$ generation symmetry---while the 5D warp factor determines the \emph{values}---the mass hierarchy and the CKM/PMNS mixing angles.


\subsection{The Dixon Algebra and the Three-Generation Problem}
\label{sec:4-dixon}

\paragraph{The problem.}
Within the CCM framework, the orbifold $S^1/\mathbb{Z}_2$ cannot fix $N_g$. The Atiyah--Patodi--Singer index theorem on $M_4 \times [0, y_c]$ with orbifold boundary conditions yields
\begin{equation}
  \mathrm{Index}(D_5) = \int_{M_4 \times I} \hat{A}(R)\, \mathrm{ch}(F) - \tfrac{1}{2}\bigl[\eta(D_{\partial}) + \dim\ker D_{\partial}\bigr].
  \label{eq:4-75}
\end{equation}
For flat $M_4$ and the topologically trivial interval $I = [0, y_c]$ (which satisfies $H^n(I, \mathbb{Z}) = 0$ for all $n \geq 1$), both the bulk integral and the boundary $\eta$-invariant vanish. There is no topological charge in the extra dimension to source a generation number. The path to fixing $N_g = 3$ therefore requires extending the \emph{finite algebra} $\mathcal{A}_F$ to one whose representation theory forces three generations.

\paragraph{The complexified Dixon algebra.}
The Dixon algebra is the tensor product of the three normed division algebras over $\mathbb{C}$~\cite{ch4:Furey2025,ch4:Dixon1994}:
\begin{equation}
  \mathcal{T}_{\mathbb{C}} = \mathbb{C}_{\mathbb{C}} \otimes_{\mathbb{C}} \mathbb{H}_{\mathbb{C}} \otimes_{\mathbb{C}} \mathbb{O}_{\mathbb{C}},
  \label{eq:4-76}
\end{equation}
where $\mathbb{C}_{\mathbb{C}} = \mathbb{C}$, $\mathbb{H}_{\mathbb{C}} = \mathbb{C} \otimes_{\mathbb{R}} \mathbb{H} \cong M_2(\mathbb{C})$, and $\mathbb{O}_{\mathbb{C}} = \mathbb{C} \otimes_{\mathbb{R}} \mathbb{O}$ is the 8-dimensional complex alternative algebra (non-associative). The complex dimensions are:
\begin{equation}
  \dim_{\mathbb{C}} \mathcal{T}_{\mathbb{C}} = 1 \times 4 \times 8 = 32.
  \label{eq:4-77}
\end{equation}
This is precisely the number of fermionic states (16~particles $+$ 16~antiparticles) in one Standard Model generation. The algebra is alternative (every subalgebra generated by two elements is associative) but not associative, owing to the octonionic factor.

\paragraph{The $\mathbb{Z}_2^5$-grading.}
The Dixon algebra $\mathcal{T} = \mathbb{C} \otimes \mathbb{H} \otimes \mathbb{O}$ carries five independent $\mathbb{Z}_2$ gradings~\cite{ch4:Furey2025}: complex conjugation on $\mathbb{C}$, the main involution on $\mathbb{H}$, octonionic conjugation on $\mathbb{O}$, and two mixed gradings from the tensor product structure. The group $\mathbb{Z}_2^5$ has $2^5 = 32$ elements, and after complexification each grading sector is one-dimensional over $\mathbb{C}$, exhausting $\mathcal{T}_{\mathbb{C}}$. The 32 grading sectors are in one-to-one correspondence with the 32 fermionic states of one SM generation: the quantum numbers (colour, weak isospin, hypercharge, chirality, particle/antiparticle) are determined by the five grading parities.

\paragraph{The KO-dimension chain.}
The Meridian geometry $M_4 \times S^1/\mathbb{Z}_2 \times F$ contributes three Clifford algebra factors:
\begin{equation}
  \mathrm{Cl}(4) \otimes \mathrm{Cl}(1) \otimes \mathrm{Cl}(6) = \mathrm{Cl}(11),
  \label{eq:4-78}
\end{equation}
with the KO-dimension sum $4 + 1 + 6 = 11$. The $\mathbb{Z}_2$ orbifold projection selects the even subalgebra:
\begin{equation}
  \mathrm{Cl}^+(11) \cong \mathrm{Cl}(10).
  \label{eq:4-79}
\end{equation}
After complexification, $\mathrm{Cl}(10)_{\mathbb{C}} = M_{32}(\mathbb{C})$, which has a \emph{unique} irreducible representation of dimension~32---matching $\dim_{\mathbb{C}} \mathcal{T}_{\mathbb{C}}$ exactly. The $\mathbb{Z}_2$ orbifold of the RS geometry naturally produces the $\mathrm{Cl}(10)$ structure within which the octonionic construction lives.


\subsection{Four Proofs of \texorpdfstring{$N_g = 3$}{Ng = 3}}
\label{sec:4-ng-proofs}

The number of fermion generations is determined by the number of independent complex structures on the octonions. A \emph{complex structure} on $\mathbb{O}$ is a linear map $J\colon \mathbb{O} \to \mathbb{O}$ (viewed as $\mathbb{R}^8$) satisfying $J^2 = -\mathrm{Id}$ and compatible with the norm: $|J(x)| = |x|$. Right multiplication by any unit imaginary octonion $e \in \mathrm{Im}(\mathbb{O})$ defines such a map: $J_e(x) = x \cdot e$, with $J_e^2 = -\mathrm{Id}$ following from alternativity and $e^2 = -1$.

While $\mathbb{O}$ has seven imaginary units $e_1, \ldots, e_7$, only three of the corresponding complex structures are independent (up to $G_2 = \mathrm{Aut}(\mathbb{O})$ equivalence after fixing the $\mathrm{SU}(3) = \mathrm{Stab}_{G_2}(e_1)$ subgroup). Each complex structure $J_i$ defines a different $\mathbb{Z}_2^5$-grading of $\mathcal{T}_{\mathbb{C}}$, and each grading yields one generation of SM fermions. The three gradings are related by the $S_3$ subgroup of $\mathrm{Aut}(\mathbb{O})$ that permutes the Fano lines through a given point.

\begin{theorem}[Algebraic Rigidity of $N_g = 3$]\label{thm:4-3}
The octonions $\mathbb{O}$ possess exactly three independent complex structures, so that $N_g = 3$ in the Dixon algebra construction. This is established by four independent algebraic results:

\begin{enumerate}
  \item \textbf{Hurwitz's theorem.} $\mathbb{R}$, $\mathbb{C}$, $\mathbb{H}$, and $\mathbb{O}$ are the only normed division algebras over $\mathbb{R}$~\cite{ch4:Hurwitz1898}. The quaternions $\mathbb{H}$ are the largest \emph{associative} normed division algebra, with $\dim_{\mathbb{R}} \mathbb{H} = 4$. The three independent complex structures on $\mathbb{O}$ arise from any quaternionic subalgebra $\mathbb{H} \hookrightarrow \mathbb{O}$: the three imaginary units of $\mathbb{H}$ define three independent $J_i$, and no fourth can be added without either lying in the $G_2$-orbit of an existing one or violating normed-division-algebra constraints.

  \item \textbf{Albert's theorem.} The exceptional Jordan algebra $J_n(\mathbb{O})$ (Hermitian $n \times n$ matrices over $\mathbb{O}$ with symmetrized product) exists only for $n \leq 3$~\cite{ch4:Albert1934,ch4:TodorovDV2018}. For $n = 3$, $J_3(\mathbb{O})$ has rank~3 (every element has exactly three eigenvalues), and $\mathrm{Aut}(J_3(\mathbb{O})) = F_4$. The maximal rank~3 bounds the number of generations that can be accommodated within an octonionic Jordan structure.

  \item \textbf{Triality uniqueness.} Among all simple Lie groups, $\mathrm{Spin}(8) = D_4$ is the \emph{unique} group possessing an outer automorphism of order~3~\cite{ch4:Baez2002}. This triality permutes three inequivalent 8-dimensional representations $\mathbf{8}_v \to \mathbf{8}_s \to \mathbf{8}_c \to \mathbf{8}_v$ and provides the $S_3$ permutation symmetry between generations. Since $\mathrm{Cl}(8) \subset \mathrm{Cl}(10)$ is naturally embedded in the Meridian Clifford chain, the triality structure is intrinsic to the framework.

  \item \textbf{Fano plane structure.} The multiplication table of $\mathrm{Im}(\mathbb{O})$ is encoded in the Fano plane (7~points, 7~lines), where each point lies on exactly 3~lines~\cite{ch4:Baez2002}. The three lines through a given point $e_i$ determine the three complex-conjugate pairs in the decomposition $\mathbb{O} \cong \mathbb{C}^4$ under $J_{e_i}$. There are exactly three independent such decompositions, corresponding to three imaginary units forming an associative triple (a single Fano line).
\end{enumerate}
\end{theorem}

\begin{proof}[Proof outline]
(1)~is classical~\cite{ch4:Hurwitz1898}. (2)~Albert showed that $J_n(\mathbb{O})$ fails to satisfy the Jordan identity for $n \geq 4$ because the non-associativity of~$\mathbb{O}$ makes $n \times n$ Hermitian matrix products ill-defined~\cite{ch4:Albert1934}. (3)~follows from inspection of the Dynkin diagrams: $D_4$ is the unique diagram with $\mathbb{Z}_3$ symmetry (and full outer automorphism group $S_3$). (4)~is verified by direct enumeration: the seven imaginary units modulo the $G_2$ action on $S^6 \subset \mathrm{Im}(\mathbb{O})$ yield $\mathrm{Stab}_{G_2}(e_1) = \mathrm{SU}(3)$, under which $\mathrm{Im}(\mathbb{O}) = \langle e_1 \rangle \oplus V_3 \oplus \bar{V}_3$, and the three independent complex structures are right multiplication by the three imaginary units of any quaternionic subalgebra $\mathbb{H} \hookrightarrow \mathbb{O}$.
\end{proof}

\begin{remark}
The three generations come from the octonionic \emph{algebra} (three complex structures on~$\mathbb{O}$), not from the Spin(10) spinor decomposition. Under $\mathrm{Spin}(10) \supset \mathrm{Spin}(8) \times \mathrm{U}(1)$, the Dirac spinor branches as $\mathbf{32} \to (\mathbf{8}_s \oplus \mathbf{8}_c)_{+1/2} \oplus (\mathbf{8}_c \oplus \mathbf{8}_s)_{-1/2}$, yielding only \emph{two} sectors. The $\mathbf{8}_v$ is absent from the spinor decomposition; it enters through the octonionic \emph{multiplication} (the algebra), not through the \emph{representation} (the spinor). Triality provides the $S_3$ symmetry that permutes the three generations---it does not generate them via the spinor branching rule.
\end{remark}


\subsection{The Dirac Operator \texorpdfstring{$D_{\mathrm{oct}}$}{D\_oct}}
\label{sec:4-doct}

The octonionic Hilbert space is $\mathcal{H}_{\mathrm{oct}} = \mathcal{H}_1 \oplus \mathcal{H}_2 \oplus \mathcal{H}_3$, where each $\mathcal{H}_i \cong \mathbb{C}^{32}$ carries the left regular representation of $\mathcal{T}_{\mathbb{C}}$ under the $i$-th $\mathbb{Z}_2^5$-grading (defined by complex structure $J_i$). The total dimension is $\dim_{\mathbb{C}} \mathcal{H}_{\mathrm{oct}} = 96$.

The finite Dirac operator $D_{\mathrm{oct}}$ is a $96 \times 96$ Hermitian matrix with block structure
\begin{equation}
  D_{\mathrm{oct}} = \begin{pmatrix}
    D_{\mathrm{gen}} & \tfrac{1}{2}\, D_{\mathrm{gen}} & \tfrac{1}{2}\, D_{\mathrm{gen}} \\[4pt]
    \tfrac{1}{2}\, D_{\mathrm{gen}} & D_{\mathrm{gen}} & \tfrac{1}{2}\, D_{\mathrm{gen}} \\[4pt]
    \tfrac{1}{2}\, D_{\mathrm{gen}} & \tfrac{1}{2}\, D_{\mathrm{gen}} & D_{\mathrm{gen}}
  \end{pmatrix},
  \label{eq:4-80}
\end{equation}
where $D_{\mathrm{gen}}$ is the standard CCM one-generation Dirac operator (a $32 \times 32$ matrix encoding the Yukawa couplings $Y_\nu$, $Y_e$, $Y_u$, $Y_d$ and the Majorana mass $M_R$ for a single generation~\cite{ch4:CCM2007,ch4:vanSuijlekom2024}). The intra-generation blocks $D_{ii} = D_{\mathrm{gen}}$ are identical for all three generations, guaranteed by the $S_3$ symmetry of the three complex structures. The inter-generation blocks $D_{ij} = \tfrac{1}{2}\, D_{\mathrm{gen}}$ ($i \neq j$) are determined by the \emph{overlap matrix} $M_{\mathrm{oct}}$.

\paragraph{The overlap matrix.}
For each complex structure $J_i$, the complexified octonions decompose as $\mathbb{O}_{\mathbb{C}} = V_i^+ \oplus V_i^-$ (eigenspaces of $J_i$ with eigenvalues $\pm i$), where $\dim_{\mathbb{C}} V_i^{\pm} = 4$. The projector onto $V_i^+$ is $P_i = \tfrac{1}{2}(\mathrm{Id} - i\, J_i)$. The mixing matrix is
\begin{equation}
  (M_{\mathrm{oct}})_{ij} = \frac{1}{\dim V^+}\, \mathrm{Tr}(P_i P_j) = \frac{1}{4}\, \mathrm{Tr}(P_i P_j).
  \label{eq:4-81}
\end{equation}

\begin{theorem}[Democratic Inter-Generation Mixing]\label{thm:4-4}
The overlap matrix is the democratic matrix
\begin{equation}
  M_{\mathrm{oct}} = \frac{1}{2}(I_3 + J_3) = \begin{pmatrix}
    1 & 1/2 & 1/2 \\ 1/2 & 1 & 1/2 \\ 1/2 & 1/2 & 1
  \end{pmatrix},
  \label{eq:4-82}
\end{equation}
where $J_3$ is the $3 \times 3$ all-ones matrix. Its eigenvalues are
\begin{equation}
  \lambda_1 = \lambda_2 = \tfrac{1}{2} \quad (\text{doubly degenerate}), \qquad \lambda_3 = 2,
  \label{eq:4-83}
\end{equation}
with eigenvectors spanning the two-dimensional irreducible and the trivial representations of~$S_3$, respectively.
\end{theorem}

\begin{proof}
The three complex structures $J_1, J_2, J_4$ (right multiplication by $e_1, e_2, e_4$) are related by the $S_3$ subgroup of $\mathrm{Aut}(\mathbb{O}) = G_2$ permuting the three Fano lines through any given point. This $S_3$ acts transitively on $\{J_1, J_2, J_4\}$, forcing: (i)~$\mathrm{Tr}(P_i P_j)$ to be the same for all pairs $i \neq j$; (ii)~$\mathrm{Tr}(P_i P_i)$ to be the same for all~$i$. The diagonal entries are $1$ by normalization ($P_i^2 = P_i$, $\mathrm{Tr}(P_i) = 4$, so $\mathrm{Tr}(P_i P_i)/4 = 1$). The off-diagonal entries are $1/2$: the singular values of the overlap matrices $U_{ij} = (V_i^+)^\dagger V_j^+$ are uniformly $1/\sqrt{2}$ for all three pairs, meaning the $+i$~eigenspaces of different complex structures overlap at exactly $45^\circ$---maximally non-aligned subject to being 4-dimensional subspaces of $\mathbb{C}^8$. The eigenvalue computation is elementary: $M_{\mathrm{oct}} = \tfrac{1}{2} I_3 + \tfrac{1}{2} J_3$, and $J_3$ has eigenvalues $\{0, 0, 3\}$.
\end{proof}

\begin{remark}
The democratic form of $M_{\mathrm{oct}}$ is forced by the $S_3$ symmetry---it has zero free parameters. The eigenvalue ratio $1\!:\!1\!:\!4$ is not the observed fermion mass hierarchy (which has ratios $\sim\!10^{-5}\!:\!10^{-2}\!:\!1$ for up-type quarks). The observed hierarchy arises from the 5D warp factor (Section~\ref{sec:4-yukawa-assessment}).
\end{remark}

\paragraph{Properties of $D_{\mathrm{oct}}$.}
The operator $D_{\mathrm{oct}}$ is Hermitian, has rank~96 (full rank, no zero modes), and its spectrum consists of 48~positive and 48~negative eigenvalues. The real structure compatibility $\|D_{\mathrm{oct}} - J_{\mathrm{oct}}\, D_{\mathrm{oct}}^*\, J_{\mathrm{oct}}^{-1}\| = 0$ holds exactly for KO-dimension~6. The grading condition $\{D_{\mathrm{oct}}, \gamma_{\mathrm{oct}}\} = 0$ is satisfied in the proper CCM basis where the chirality decomposition ($\gamma_{\mathrm{oct}} = \mathrm{diag}(+I_{48}, -I_{48})$) aligns with the particle--antiparticle decomposition: since $D_{\mathrm{gen}}$ is off-diagonal in the chirality basis (the Yukawa matrix connects $L$ to $R$), the inter-generation mixing $M_{\mathrm{oct}} \otimes D_{\mathrm{gen}}$ inherits this off-diagonal structure.


\subsection{Poincar\'{e} Duality via the Associative Envelope}
\label{sec:4-poincare}

The Poincar\'{e} duality axiom requires non-degeneracy of the intersection form on $K_*(\mathcal{A}_F) \times K_*(\mathcal{A}_F) \to \mathbb{Z}$. For the associative CCM algebra $\mathcal{A}_F = \mathbb{C} \oplus \mathbb{H} \oplus M_3(\mathbb{C})$, this follows from the standard K-theory of direct sums of matrix algebras: $K_0(\mathcal{A}_F) = \mathbb{Z}^3$ and the intersection form is a non-degenerate $3 \times 3$ integer matrix~\cite{ch4:CCM2007}. For the non-associative Dixon algebra $\mathcal{T}_{\mathbb{C}}$, standard K-theory does not directly apply: matrix algebras $M_n(\mathbb{O}_{\mathbb{C}})$ are ill-defined for $n > 1$ because non-associativity prevents consistent matrix multiplication.

The resolution passes through the \emph{associative envelope} of the octonions.

\begin{theorem}[Associative Envelope and Poincar\'{e} Duality]\label{thm:4-5}
\leavevmode
\begin{enumerate}
  \item[\emph{(a)}] The associative envelope of the octonions is $A_{\mathrm{env}}(\mathbb{O}) \cong M_8(\mathbb{R})$.
  \item[\emph{(b)}] The associative envelope of the full Dixon algebra is $A_{\mathrm{env}}(\mathcal{T}_{\mathbb{C}}) \cong M_{16}(\mathbb{C})$.
  \item[\emph{(c)}] $K_0(A_{\mathrm{env}}(\mathcal{T}_{\mathbb{C}})) = \mathbb{Z}$ and $K_1(A_{\mathrm{env}}(\mathcal{T}_{\mathbb{C}})) = 0$ (Morita equivalence $M_{16}(\mathbb{C}) \sim \mathbb{C}$).
  \item[\emph{(d)}] The representation K-theory $K_0(\pi(\mathcal{A}_{\mathrm{oct}}) \text{ on } \mathcal{H}_{\mathrm{oct}})$ is isomorphic to~$\mathbb{Z}^3$, matching the CCM result.
  \item[\emph{(e)}] The intersection form is non-degenerate, and Poincar\'{e} duality holds.
\end{enumerate}
\end{theorem}

\begin{proof}
For~(a): the envelope is generated by the left and right multiplication operators $\{L_{e_i}, R_{e_i} : i = 0, \ldots, 7\}$, which are elements of $\mathrm{End}_{\mathbb{R}}(\mathbb{O}) = M_8(\mathbb{R})$. The linear span of $\{L_i, R_i\}$ alone has dimension~15 ($\mathrm{Id}$ plus 7~left and 7~right operators). Including pairwise products $\{L_i L_j, L_i R_j, R_i L_j, R_i R_j\}$ increases the dimension to $64 = \dim M_8(\mathbb{R})$, so the envelope is all of $M_8(\mathbb{R})$.

For~(b): the complexified factors satisfy $(\mathbb{C}_{\mathbb{C}})_{\mathrm{env}} = \mathbb{C}$ and $(\mathbb{H}_{\mathbb{C}})_{\mathrm{env}} = M_2(\mathbb{C})$ (both already associative), and $A_{\mathrm{env}}(\mathbb{O}_{\mathbb{C}}) = M_8(\mathbb{C})$ (complexification of part~(a)). Tensoring gives $\mathbb{C} \otimes M_2(\mathbb{C}) \otimes M_8(\mathbb{C}) = M_{16}(\mathbb{C})$.

For~(c): $M_{16}(\mathbb{C})$ is Morita equivalent to $\mathbb{C}$, giving $K_0 = \mathbb{Z}$ and $K_1 = 0$.

For~(d): the gauge group extraction (Section~\ref{sec:4-gauge-oct}) gives $G_{\mathrm{SM}} = \mathrm{SU}(3) \times \mathrm{SU}(2) \times \mathrm{U}(1)$, unchanged from CCM. The representation of $\mathcal{A}_{\mathrm{oct}}$ on $\mathcal{H}_{\mathrm{oct}} = \mathbb{C}^{96}$ decomposes into the same irreducible pieces as in CCM (three generations of $16 + 16$ states with identical gauge quantum numbers), so the K-theoretic data is isomorphic.

Part~(e) follows from~(d) and the explicit CCM intersection form computation~\cite{ch4:CCM2007}.
\end{proof}

\begin{remark}[K-theoretic stability]
The octonionic extension is a ``soft deformation'' in the K-theoretic sense. The non-associativity introduces corrections proportional to the octonionic associator $[a, b, c] = (ab)c - a(bc)$, but these corrections: (i)~preserve representation dimensions ($\dim \mathcal{H}_{\mathrm{oct}} = 96$); (ii)~preserve the gauge group; (iii)~vanish on the associative subalgebra generating the gauge group; and (iv)~do not change the homotopy type of the algebra. The K-groups are therefore stable under the non-associative deformation.
\end{remark}

\begin{remark}[Structural necessity of the envelope]\label{rem:4-envelope-necessity}
One may ask whether the associative envelope detour in Theorem~\ref{thm:4-5} can be avoided entirely---whether $K_0(\mathbb{O}) = \mathbb{Z}$ can be proven using only the alternative identities and Artin's theorem. We investigated this in detail, and the answer is \emph{no}: the envelope is structurally necessary.

The obstruction is precise. The standard proof that every finitely generated module over a division algebra is free requires the splitting step: if $a \cdot m_1 = b \cdot m_2$ with $b \neq 0$, then $m_2 = b^{-1} \cdot (a \cdot m_1) = (b^{-1} a) \cdot m_1$. This last equality requires $[b^{-1}, a, m_1] = 0$---a \emph{three-element} expression. Artin's theorem guarantees associativity only for subalgebras generated by \emph{two} elements, and numerical verification confirms that $[b^{-1}, a, m]$ is non-zero for $100\%$ of random octonionic triples. The boundary is tight: all two-element module expressions (left alternative, right alternative) hold exactly, while all three-element expressions fail generically.

Consequently, the K-group computation $K_0(\mathbb{O}) = \mathbb{Z}$ requires the associative envelope. This is not a proof gap but a structural feature: K-theory depends on module categories, module decomposition requires three-element associativity, and three-element associativity requires the envelope. For any alternative algebra~$A$, we have $K_0^{\mathrm{alt}}(A) \cong K_0(A_{\mathrm{env}})$ because every $A$-module action factors through $A \to \mathrm{End}(A) \cong A_{\mathrm{env}}$.
\end{remark}

\begin{remark}[Direct intersection form]\label{rem:4-direct-intersection}
While the K-group computation requires the envelope, the \emph{intersection form} can be constructed directly from octonionic data. The three complex structures $J_a = L_{e_a}$ (left multiplication by imaginary octonions $e_1, e_2, e_4$) define projectors $P_a = \frac{1}{2}(\mathrm{Id} - i J_a)$ whose overlaps give
\begin{equation}
  M_{\mathrm{oct}}(a,b) = \frac{\mathrm{Tr}(P_a P_b)}{\mathrm{Tr}(P_a)} = \frac{8 - \mathrm{Tr}(J_a J_b)}{16} = \begin{cases} 1 & a = b \\ 1/2 & a \neq b \end{cases},
  \label{eq:4-intersection-direct}
\end{equation}
reproducing the democratic matrix of Theorem~\ref{thm:4-4}. The off-diagonal value $1/2$ follows from the anticommutativity of distinct imaginary unit operators: $L_{e_a} L_{e_b} = -L_{e_b} L_{e_a}$ for $a \neq b$ implies $\mathrm{Tr}(L_{e_a} L_{e_b}) = 0$ by the cyclic property of the trace. The determinant $\det M_{\mathrm{oct}} = 1/2 \neq 0$ establishes non-degeneracy over~$\mathbb{Q}$, confirming Poincar\'{e} duality without invoking the envelope. No associative scaffolding was used in this construction---only the octonionic multiplication table and the complex structure overlap formula.
\end{remark}


\subsection{Orientability and the \texorpdfstring{$1/6$}{1/6} Coefficient}
\label{sec:4-orientability-oct}

The orientability axiom for a spectral triple $(\mathcal{A}, \mathcal{H}, D, J, \gamma)$ requires the existence of a Hochschild $n$-cycle $c \in Z_n(\mathcal{A}, \mathcal{A} \otimes \mathcal{A}^{\mathrm{op}})$ with $\pi(c) = \gamma$, where $\gamma$ is the grading and $n$ is the metric dimension~\cite{ch4:Connes1994}. For the finite spectral triple (metric dimension~0), the Hochschild cycle is a 0-cycle---an element $c \in \mathcal{A} \otimes \mathcal{A}^{\mathrm{op}}$ satisfying the trivial condition $d_0(c) = 0$ (since $d_0$ maps to the trivial degree~$-1$ module). The axiom thus reduces to a \emph{surjectivity condition}: does $\gamma$ lie in the image of the bimodule representation $\pi\colon \mathcal{A} \otimes \mathcal{A}^{\mathrm{op}} \to \mathrm{End}(\mathcal{H})$?

For non-associative algebras, three complications arise: (i)~the standard Hochschild differential satisfies $d^2 \neq 0$ at degrees $\geq 2$; (ii)~the nucleus $N(\mathbb{O}) = \{n \in \mathbb{O} : [n, a, b] = 0\ \forall\, a, b\} = \mathbb{R} \cdot e_0$ is too small to contain the grading; (iii)~the bimodule structure involves parenthesisation. However, at degree~0 none of these obstructions apply: $d_0 = 0$ regardless of associativity, and the surjectivity of the bimodule map can be established directly.

\begin{theorem}[Orientability of the Octonionic Spectral Triple]\label{thm:4-6}
The orientability axiom holds for $(\mathcal{T}_{\mathbb{C}}, \mathcal{H}_{\mathrm{oct}}, D_{\mathrm{oct}}, J_{\mathrm{oct}}, \gamma_{\mathrm{oct}})$, established by two independent proofs.

\textbf{Proof~A (Associative envelope surjectivity).}
From Theorem~\ref{thm:4-5}(a), $A_{\mathrm{env}}(\mathbb{O}) = M_8(\mathbb{R})$. The bimodule map $\mu\colon \mathbb{O} \otimes_{\mathbb{R}} \mathbb{O}^{\mathrm{op}} \to \mathrm{End}_{\mathbb{R}}(\mathbb{O})$, defined by $\mu(a \otimes b^{\mathrm{op}})(x) = (a \cdot x) \cdot b$, has image $\mathrm{span}\{L_{e_i} R_{e_j} : i, j = 0, \ldots, 7\}$. This span has rank~$64 = \dim M_8(\mathbb{R})$, so $\mu$ is surjective. In particular, the grading $\gamma_{\mathbb{O}} = \mathrm{diag}(+1, -1, \ldots, -1)$ (octonionic conjugation: $\gamma_{\mathbb{O}}(x) = x^*$) lies in $\mathrm{Im}(\mu)$. For the full algebra $A_{\mathrm{env}}(\mathcal{T}_{\mathbb{C}}) = M_{16}(\mathbb{C})$, surjectivity of the bimodule map follows from the double commutant theorem (Burnside's theorem for simple algebras). The three-generation grading $\gamma_{\mathrm{oct}} = I_3 \otimes \gamma_{1\mathrm{gen}}$ lies in the image.

\textbf{Proof~B (Explicit Hochschild cycle).}
For a normed division algebra $K$ of real dimension~$n$, the universal trace formula~\cite{ch4:Baez2002} gives
\begin{equation}
  \sum_{i=0}^{n-1} (e_i \cdot x) \cdot e_i = (2 - n)\, x^*,
  \label{eq:4-84}
\end{equation}
where $\{e_0, \ldots, e_{n-1}\}$ is an orthonormal basis and $x^*$ is the conjugation. For $\mathbb{O}$ ($n = 8$):
\begin{equation}
  \sum_{i=0}^{7} (e_i \cdot x) \cdot e_i = -6\, x^*.
  \label{eq:4-85}
\end{equation}
Define the orientation cycle
\begin{equation}
  c_{\mathbb{O}} = -\frac{1}{6}\, \sum_{i=0}^{7} e_i \otimes e_i^{\mathrm{op}} \;\in\; \mathbb{O} \otimes_{\mathbb{R}} \mathbb{O}^{\mathrm{op}}.
  \label{eq:4-86}
\end{equation}
Then
\begin{equation}
  \pi(c_{\mathbb{O}})(x) = -\frac{1}{6}\, \sum_{i=0}^{7} (e_i \cdot x) \cdot e_i = -\frac{1}{6}(-6\, x^*) = x^* = \gamma_{\mathbb{O}}(x).
  \label{eq:4-87}
\end{equation}
The element $c_{\mathbb{O}}$ is a Hochschild 0-cycle trivially ($d_0 = 0$), and the full orientation cycle is $c_{\mathrm{oct}} = I_3 \otimes c_{\mathbb{C}} \otimes c_{\mathbb{H}} \otimes c_{\mathbb{O}}$, where $c_{\mathbb{C}}$ and $c_{\mathbb{H}}$ are the standard orientation cycles for the associative factors.
\end{theorem}

\begin{remark}
The coefficient $-1/6$ in~\eqref{eq:4-86} arises from $1/(n - 2) = 1/(8 - 2) = 1/6$, where $n = 8$ is the octonionic dimension. This is numerically equal to the conformal coupling $\xi = 1/6$ derived in Section~\ref{sec:4-xi-derivation}, which arises from $\xi = (d - 2)/[4(d - 1)] = 2/12 = 1/6$ for spacetime dimension $d = 4$. Both are consequences of dimensionality (8~for octonions, 4~for spacetime) through different mechanisms. Whether this coincidence has deeper significance through the Meridian bulk--brane correspondence---where the 8-dimensional octonionic structure on the brane is related to 4-dimensional spacetime through the RS warped geometry---is an open question noted here as an observation.
\end{remark}

\paragraph{Basis independence.}
The orientation cycle $c_{\mathbb{O}}$ is basis-independent: for any orthonormal basis $\{f_0, \ldots, f_7\}$ of $\mathbb{O}$ with $f_0 = e_0$, the trace operator $T(x) = \sum_i f_i\, x\, f_i$ is invariant under $\mathrm{Aut}(\mathbb{O}) = G_2$, which acts transitively on orthonormal bases of $\mathrm{Im}(\mathbb{O})$.


\subsection{Complete Axiom Verification}
\label{sec:4-axiom-table}

With orientability established, the octonionic spectral triple satisfies all seven NCG axioms:

\begin{table}[htbp]
\centering
\caption{NCG axiom verification for the octonionic spectral triple $(\mathcal{T}_{\mathbb{C}}, \mathcal{H}_{\mathrm{oct}}, D_{\mathrm{oct}}, J_{\mathrm{oct}}, \gamma_{\mathrm{oct}})$.}
\label{tab:4-4}
\small
\begin{tabular}{@{}clll@{}}
\toprule
\# & Axiom & Status & Key technique \\
\midrule
1 & Compact resolvent & Holds & Finite-dimensional (trivial) \\
2 & First-order condition & Modified$^\dagger$ & Boyle--Farnsworth~\cite{ch4:BoyleFarnsworth2020}: holds up to associator \\
3 & Orientability & Holds & Trace formula $+$ envelope surjectivity (Theorem~\ref{thm:4-6}) \\
4 & Poincar\'{e} duality & Holds & Via $A_{\mathrm{env}} = M_8(\mathbb{R})$, $K_0 = \mathbb{Z}$ (Thm.~\ref{thm:4-5}); direct form~\eqref{eq:4-intersection-direct} \\
5 & Reality & Holds & Conjugations on $\mathbb{C}$, $\mathbb{H}$, $\mathbb{O}$ are anti-involutions \\
6 & Finiteness & Holds & Finite-dimensional (trivial) \\
7 & Regularity & Holds & Finite-dimensional (trivial) \\
\bottomrule
\multicolumn{4}{@{}l}{\footnotesize $^\dagger$\,The first-order condition $[[D, a], Jb^* J^{-1}] = 0$ holds \emph{exactly} on the}\\
\multicolumn{4}{@{}l}{\footnotesize \phantom{$^\dagger$\,}associative subalgebra $\mathbb{C} \oplus \mathbb{H} \oplus M_3(\mathbb{C})$ (which generates $G_{\mathrm{SM}}$). On the full}\\
\multicolumn{4}{@{}l}{\footnotesize \phantom{$^\dagger$\,}octonionic part, the violation is bounded by the associator: $[[D, a], Jb^* J^{-1}] = \Delta(a, b)$,}\\
\multicolumn{4}{@{}l}{\footnotesize \phantom{$^\dagger$\,}where $\Delta$ is totally antisymmetric and vanishes on the physical gauge sector.}\\
\end{tabular}
\end{table}


\subsection{The Yukawa Structure: An Honest Assessment}
\label{sec:4-yukawa-assessment}

The octonionic associator $[a, b, c] = (ab)c - a(bc)$ is totally antisymmetric (a 3-form on $\mathrm{Im}(\mathbb{O})$), vanishes on the 7~Fano triples (associative sub-triples), and is non-zero on exactly 28 of the 35 independent triples of imaginary units~\cite{ch4:Baez2002}. The 35-parameter space $\Lambda^3(\mathbb{R}^7)$ decomposes under $G_2 = \mathrm{Aut}(\mathbb{O})$ as
\begin{equation}
  \Lambda^3(\mathbb{R}^7) = \mathbf{1} \oplus \mathbf{7} \oplus \mathbf{27},
  \label{eq:4-88}
\end{equation}
where the $\mathbf{1}$ is the $G_2$-invariant associative calibration $\varphi$ (non-zero on Fano triples, zero elsewhere) and the associator itself lies in the $\mathbf{7} \oplus \mathbf{27}$ complement.

A critical point: \emph{the octonionic associator has zero free parameters}. It is completely determined by the Fano plane. The 35 independent components are all fixed---28 non-zero values of magnitude~$2$ with determined signs, and 7~zeros on the Fano lines. Therefore the association between ``35~parameters'' and the 20~physical SM mixing parameters (4~CKM, 4~PMNS, 12~fermion masses) is a counting coincidence, not a physical mechanism.

\paragraph{The $S_3$ obstruction.}
The three complex structures $J_1, J_2, J_4$ are related by an exact $S_3$ symmetry, forcing $M_{\mathrm{oct}}$ to be democratic (Theorem~\ref{thm:4-4}). The up-type and down-type quarks see the \emph{same} overlap matrix, so the CKM matrix at the algebraic level is
\begin{equation}
  V_{\mathrm{CKM}}^{\mathrm{alg}} = U_u^\dagger\, U_d = U_M^\dagger\, U_M = I_3,
  \label{eq:4-89}
\end{equation}
where $U_M$ diagonalises $M_{\mathrm{oct}}$. This predicts zero CKM mixing---obviously wrong.

\paragraph{Bulk--brane complementarity.}
The resolution is the 5D warp factor. The effective Yukawa coupling between generations $i$ and $j$ for fermion type~$f$ is
\begin{equation}
  Y_{ij}^{f,\,\mathrm{eff}} = (M_{\mathrm{oct}})_{ij}\, y_f \int_0^{y_c} f_i^f(y)\, f_j^f(y)\, h(y)\, e^{-4A(y)}\, dy,
  \label{eq:4-90}
\end{equation}
where $f_i^f(y)$ is the extra-dimensional profile of the $i$-th generation fermion of type~$f$, $h(y)$ is the Higgs profile, and $A(y) = -ky$ is the warp factor. The profiles depend on the bulk mass parameters~$c_i^f$, which are different for up-type and down-type quarks. The exponential $e^{-c\, k y_c}$ with $ky_c \approx 37$ converts $\order{1}$ differences in~$c$ to enormous mass hierarchies ($e^{-0.5 \times 37} \sim 10^{-8}$, comparable to $m_u/m_t$), and the difference between $Y_u^{\mathrm{eff}}$ and $Y_d^{\mathrm{eff}}$ generates non-trivial CKM mixing.

\begin{table}[htbp]
\centering
\caption{Division of labour between the algebraic and geometric sectors.}
\label{tab:4-5}
\small
\begin{tabular}{@{}lll@{}}
\toprule
Structure & Determines & Free parameters \\
\midrule
Octonionic (brane) & $N_g = 3$, gauge group, $S_3$ symmetry, Fano topology & 0 \\
5D warp (bulk) & Mass hierarchy, CKM/PMNS mixing, flavour violation & $\sim$20 (bulk masses $c_i^f$) \\
\bottomrule
\end{tabular}
\end{table}

This is not a weakness but the correct architectural division: the octonions fix what \emph{can} exist (three generations with $S_3$ symmetry and identical gauge quantum numbers), while the 5D geometry fixes what \emph{does} exist in our universe (the specific mass ratios and mixing angles).


\subsection{The Gauge Group from Octonions}
\label{sec:4-gauge-oct}

The gauge group is extracted from the inner automorphisms of the Dixon algebra modulo its centre. The automorphism group of $\mathbb{O}$ is $G_2$~\cite{ch4:Baez2002}, but the choice of a preferred complex structure (equivalently, a preferred imaginary unit $e_1$) breaks $G_2$ to the stabiliser:
\begin{equation}
  \mathrm{Stab}_{G_2}(e_1) = \mathrm{SU}(3).
  \label{eq:4-91}
\end{equation}
This is the colour group. The breaking $G_2 \to \mathrm{SU}(3)$ is automatic once the generation structure is defined---it is not an additional assumption. Combined with the automorphisms of the other Dixon factors:
\begin{equation}
  G_{\mathrm{SM}} = \mathrm{Stab}_{G_2}(e_1) \times \mathrm{Aut}(\mathbb{H}) \times \mathrm{Aut}(\mathbb{C}) = \mathrm{SU}(3) \times \mathrm{SU}(2) \times \mathrm{U}(1),
  \label{eq:4-92}
\end{equation}
matching the CCM gauge group exactly~\cite{ch4:Krasnov2025,ch4:Furey2025}.

\paragraph{Spectral action preservation.}
The gravitational spectral action coefficients depend only on the 5D Dirac operator (the bulk geometry), not on the finite space~$F$. Since the octonionic extension modifies only the brane algebra while $\dim \mathcal{H}_{\mathrm{oct}} = 96 = \dim \mathcal{H}_F^{\mathrm{CCM}}$:
\begin{equation}
  (C^2,\; E_4,\; R^2) = (-18,\; +11,\; 0) \qquad \text{(unchanged)}.
  \label{eq:4-93}
\end{equation}
The $R^2 = 0$ identity (the structural relation $15/12 - 8/12 - 7/12 = 0$) is a property of the 5D Dirac operator and holds for \emph{any} finite space. The conformal coupling $\xi = 1/6$ (Section~\ref{sec:4-xi-derivation}), the self-tuning mechanism, and the sound speed $c_s \sim 10c$ are all bulk-determined and therefore preserved.

\paragraph{The matter spectral action.}
The matter sector $S_{\mathrm{matter}} = \Tr_{\mathcal{H}_F}[f(D_F^2/\Lambda^2)]$ has its leading coefficient
\begin{equation}
  a_0^{\mathrm{matter}} = \Tr(\mathbf{1}_{\mathcal{H}_F}) = 96,
  \label{eq:4-94}
\end{equation}
which is identical in both the CCM ($N_g = 3$, assumed) and octonionic ($N_g = 3$, derived) constructions. The higher Seeley--DeWitt coefficients $a_2^{\mathrm{matter}}$ and $a_4^{\mathrm{matter}}$ depend on the specific form of $D_F$ (the Yukawa matrix). If the inter-generation mixing from $M_{\mathrm{oct}}$ differs from the CCM Yukawa matrix, these coefficients will receive $\order{\tfrac{1}{2}}$ corrections from the off-diagonal blocks. The gravitational predictions of Meridian---which depend only on the bulk spectral action and the leading matter trace---are unaffected.


%% ============================================================
%% BIBLIOGRAPHY ENTRIES TO ADD to the existing \begin{thebibliography}
%% (These should be merged into the main bibliography)
%% ============================================================

%% \bibitem{ch4:BoyleFarnsworth2020}
%% L.~Boyle and S.~Farnsworth,
%% ``Non-Commutative Geometry, Non-Associative Geometry, and the Standard Model of Particle Physics,''
%% New J.\ Phys.\ \textbf{16}, 123027 (2020);
%% arXiv:1910.11888.
%%
%% \bibitem{ch4:Furey2025}
%% C.~Furey,
%% ``Standard Model Physics from an Algebra?''
%% Annalen der Physik (2025).
%%
%% \bibitem{ch4:Dixon1994}
%% G.~M.~Dixon,
%% ``Division Algebras: Octonions, Quaternions, Complex Numbers and the Algebraic Design of Physics,''
%% Springer (1994).
%%
%% \bibitem{ch4:Hurwitz1898}
%% A.~Hurwitz,
%% ``\"Uber die Komposition der quadratischen Formen von beliebig vielen Variablen,''
%% Nachr.\ Ges.\ Wiss.\ G\"ottingen (1898), 309--316.
%%
%% \bibitem{ch4:Albert1934}
%% A.~A.~Albert,
%% ``On a certain algebra of quantum mechanics,''
%% Ann.\ Math.\ \textbf{35}, 65--73 (1934).
%%
%% \bibitem{ch4:TodorovDV2018}
%% I.~Todorov and M.~Dubois-Violette,
%% ``Exceptional Quantum Geometry and Particle Physics,''
%% arXiv:1806.09450 (2018).
%%
%% \bibitem{ch4:Baez2002}
%% J.~C.~Baez,
%% ``The Octonions,''
%% Bull.\ AMS \textbf{39}, 145--205 (2002);
%% math/0105155.
%%
%% \bibitem{ch4:Krasnov2025}
%% K.~Krasnov,
%% ``The Standard Model Gauge Group from Octonionic Pure Spinors,''
%% arXiv:2504.16465 (2025).
%%
%% \bibitem{ch4:vanSuijlekom2024}
%% W.~D.~van Suijlekom,
%% ``Noncommutative Geometry and Particle Physics,''
%% 2nd edition, Springer (2024).
